\chapter{Preliminary Results and Interpretation}
\label{ch:preliminary}
\section{Completed Evidence}
The following observations are completed repository results, not outcomes
promised for the proposed study. The archived evaluation completed
\CompletedStages\ stages with \FailedStages\ stage failures and
\FailedIntegrations\ failed integrations. Numerical execution success is
separate from satisfying the accuracy criteria in \cref{tab:criteria}.
\Cref{tab:summary} summarizes the main measurements.

\begin{table}[!htbp]
\centering\small
\caption{Automatically extracted preliminary numerical evidence.}
\label{tab:summary}
\begin{tabularx}{\textwidth}{@{}p{.29\textwidth}Yr@{}}
\toprule Comparison & Metric & Value\\\midrule
\TableRows{summary}
\bottomrule
\end{tabularx}
\end{table}

\section{Analytical and Reference Checks}
\Cref{fig:poiseuille} compares analytical, FEM, and PINN straight-channel
velocity profiles. The PINN velocity error is
\PoiseuilleVelocityPercent\%, while its longitudinally sampled pressure
error is \PoiseuillePressurePercent\%. The FEM velocity error is near
floating-point/iteration accuracy for this polynomial solution. This
verifies the controlled case, not the accuracy of every curved geometry.

\ResultFigure{16_poiseuille_validation}{Poiseuille verification}{Analytical
plane Poiseuille profile with independent FEM and separately trained PINN
profiles. The second panel shows absolute velocity error. Relative pressure
errors are evaluated on longitudinally distributed points, not this section.}{fig:poiseuille}

\begin{table}[!htbp]
\centering\small
\caption{Three reference meshes and their measured sensitivity.}
\label{tab:mesh}
\begin{tabularx}{\textwidth}{@{}lrrY@{}}
\toprule Grid & Triangles & Velocity difference (\%) & $\lVert\nabla\cdot\vect u\rVert_{L^2}$\\\midrule
\TableRows{mesh}
\bottomrule
\end{tabularx}
\end{table}
The penultimate-to-finest velocity difference is
\MeshDifferencePercent\%, and the trajectory interpolator differs from
direct FEM velocity evaluation by \InterpolationPercent\%.
The zero difference in the finest row is a self-comparison; it is not an
estimate of zero discretization error. The declining divergence norm in
\cref{tab:mesh} and the sensitivity plot in \cref{fig:mesh} support reference
refinement, while leaving systematic modeling error unmeasured.

\ResultFigure{08_mesh_sensitivity}{Reference mesh sensitivity}{Velocity
differences to the finest available mesh and the integrated FEM divergence
norm. The plotted floor for the finest self-comparison is a display choice,
not a formal uncertainty estimate.}{fig:mesh}

\section{Curved-Channel Flow}
The principal seed gives sampled velocity and pressure errors of
\FlowVelocityPercent\% and \FlowPressurePercent\%, respectively.
\Cref{fig:fields} compares the spatial fields and errors;
\cref{fig:parity} shows pointwise held-out agreement. Plotting-grid reference
values use the measured fast interpolation; numerical validation metrics
use direct FEM evaluation. Maximum section-flux deviation is
\FluxPercent\%, which is below the present 5\% screen but remains a
nonzero transport inconsistency.

\WideResultFigure{04_flow_fields}{Channel fields and errors}{PINN and FEM
speed and gauge-aligned pressure with local error maps. Both methods solve
the same stated boundary-value problem. Smooth appearance does not imply
uniformly small boundary or trajectory error.}{fig:fields}
\WideResultFigure{05_validation_parity}{Field parity}{Direct FEM/PINN parity
on independent interior sample points. Pressure samples are mean-centered
before comparison.}{fig:parity}

\begin{table}[!htbp]
\centering\small
\caption{Flow error across the two recorded training seeds.}
\label{tab:seeds}
\begin{tabularx}{\textwidth}{@{}lYY@{}}
\toprule Seed & Velocity error (\%) & Pressure error (\%)\\\midrule
\TableRows{seeds}
\bottomrule
\end{tabularx}
\end{table}
The velocity error has mean \VelocityMeanPercent\% and sample standard
deviation \VelocitySDPercent\ percentage points across these two seeds.
The low observed spread does not establish a general training-success
probability; the sample is too small for that conclusion.

\section{Particle Paths and Endpoint Decisions}
\Cref{fig:particles} compares the baseline particle paths. Endpoint classes
agree for \BaselineAgreement\ of \BaselineCount\ inlet seeds. The
reference predicts exit for all seven; the surrogate predicts sticking for
one. In addition, the largest common-time RMS path discrepancy is
\MaxPathWidthPercent\% of the channel width. Thus, a field error near
1.5\% does not make the baseline particle prediction uniformly accurate.

\ResultFigure[1.0]{10_particle_paths}{Baseline particle paths}{Baseline
particle trajectories for PINN plus ForceNet and FEM plus exact Stokes drag.
Crosses mark sticking; other endpoint types are recorded in the data. The
outer-seed endpoint discrepancy is retained.}{fig:particles}

\begin{table}[!htbp]
\centering\small
\caption{Baseline particle comparisons on the common physical time interval.}
\label{tab:paths}
\begin{tabularx}{\textwidth}{@{}rYYY@{}}
\toprule Inlet $\eta$ & RMS / width (\%) & Time coverage (\%) & Same endpoint\\\midrule
\TableRows{trajectories}
\bottomrule
\end{tabularx}
\end{table}
The lowest coverage is \MinimumTimeCoveragePercent\%, so the matched-time
error alone would conceal an appreciable difference in termination time.
\Cref{tab:paths} motivates retaining all three trajectory criteria.
Comparing ForceNet with exact drag on the same PINN flow preserves all seven
baseline endpoint classes. That isolation result suggests the flow
approximation is a more consequential contributor to the observed class
discrepancy in this case, although it does not prove a single local cause.

\section{Diameter and Density Sensitivity}
The path comparisons in \cref{fig:diameter,fig:density} show the effect of the
selected inertia and size changes. \Cref{tab:outcomes} gives the actual
counts. The largest diameter sticks for every seed in both approaches;
the higher-density case retains a disagreement in the number exiting.
These observations are conditional on the drag-only model and selected
seven starting positions. They are not estimates of a device's population
deposition efficiency.

\begin{table}[!htbp]
\centering\small
\caption{Endpoint counts in the completed parameter cases; seven seeds per row.}
\label{tab:outcomes}
\begin{tabularx}{\textwidth}{@{}lYrrr@{}}
\toprule Case & Method & Exit & Stick & Other\\\midrule
\TableRows{outcomes}
\bottomrule
\end{tabularx}
\end{table}
\WideResultFigure{11_diameter_study}{Diameter sensitivity}{Particle paths
for the three recorded diameters. The largest diameter is 15\% of the width,
so physical confinement effects need further consideration.}{fig:diameter}
\WideResultFigure{12_density_study}{Density sensitivity}{Particle paths
for the three recorded density ratios under the same stated flow and
contact model.}{fig:density}

\section{Three-Dimensional Evidence}
The separate ABC model has velocity error \ABCVelocityPercent\% and
pressure error \ABCPressurePercent\%. The latter exceeds the current
5\% criterion. \Cref{fig:abc,fig:abcvalidation} demonstrates nonplanar
trajectory integration and records the associated approximation quality;
it does not provide a validated 3D curved-channel solution.

The unresolved path accuracy, endpoint agreement and 3D pressure motivate
the next experiments in \cref{ch:plan}. Additional diagnostics, including
the planar extrusion labeled as such, appear in \cref{app:figures}.

\ResultFigure[1.0]{19_3d_benchmark}{Three-dimensional tracer paths}{Tracer
paths from the 3D PINN and exact forced ABC velocity fields. These are
genuine three-coordinate paths in a separate benchmark.}{fig:abc}
\ResultFigure{20_3d_validation}{Three-dimensional benchmark checks}{Speed
parity and physics loss for the ABC benchmark. The pressure error is reported
numerically and remains above the current screening criterion.}{fig:abcvalidation}
